

Our work sits between three lines of research — agent memory systems, long-term
conversation benchmarks, and retrieval-augmented generation — and departs from
all three on one axis: we constrain the memory manager to a small *local* model
on consumer hardware and evaluate length-invariance directly.

\section{Agent memory systems}

\textbf{MemGPT} \citep{packer2023memgpt} frames the LLM as an operating system that
pages information between a bounded context window and external storage,
managing memory tiers to give the appearance of unbounded context. \textbf{Mem0}
\citep{chhikara2025mem0} builds a memory layer that extracts, consolidates, and
retrieves salient facts from ongoing conversations, reporting large token-cost
and latency savings — the closest system in spirit to our typed extraction
protocol. \textbf{Generative Agents} \citep{park2023generative} store a natural-language
memory stream retrieved by a recency/importance/relevance score and periodically
synthesised into higher-level reflections; our age-aware retrieval and the
consolidation pass we propose as future work are related in mechanism. Across
these systems the memory manager is assumed to be a strong, typically
API-hosted model. We ask the opposite question — how small can that manager be
— and quantify the answer (§5.5).

\section{Long-term conversation benchmarks}

\textbf{LoCoMo} \citep{maharana2024locomo}, our primary benchmark, provides very long
multi-session dialogues (hundreds of turns over many sessions) with QA spanning
single-hop, multi-hop, temporal, open-domain, and adversarial categories; the
original work shows that both long-context models and RAG lag human performance,
especially on temporal reasoning — the gap our grounding mechanism targets. We
add abstain-aware and LLM-judged scoring and, critically, a
length-invariance split the original evaluation does not report. **Beyond
Goldfish Memory / Multi-Session Chat** \citep{xu2022msc} studies persona
consistency across sessions and is our planned second benchmark.

\section{Retrieval-augmented generation}

\textbf{RAG} \citep{lewis2020rag} established retrieve-then-generate over a dense index
as a way to extend a model's effective knowledge; our raw-turn baseline is a
direct instance. \textbf{Self-RAG} \citep{asai2023selfrag} adds self-reflective control
over *when* to retrieve and whether generations are supported — a decision our
agentic protocol makes at *write* time (what is worth storing, whether a claim
should update or delete an existing memory) rather than at read time.

\section{Small-model capability}

A separate line examines the capability gap between small open models and
frontier APIs, and the effect of quantisation on that gap. This grounds our
framing question — *can a 3B–14B model perform reliable memory management?* —
which §5.5 answers with a floor of roughly 7–14B, with temporal grounding the
first capability to degrade under compression.

\section{Positioning}

Unlike prior memory systems, which assume a capable (often API) manager and
report aggregate accuracy at a fixed conversation length, we (i) constrain the
manager to a small local model and measure the minimum viable size, and (ii)
evaluate length-invariance directly, showing where retrieval-based memory
overtakes context-stuffing as conversations grow.